\section{经典 Floer 理论}\label{sec:8}
\subsection{Riemann 曲面}\label{sec:8a}
从这里开始，固定一个带角正合辛流形 $M$，在以下整个讨论中把它作为“目标空间”。将 $M$ 上给定的结构记为 $\varpi=\varpi_M$、$\theta=\theta_M$、$I=I_M$。
\begin{definition*}
设 $\widehat S$ 是紧致带边 Riemann 曲面，$\Sigma$ 是有限个边界点所成的集合，并分成两部分 $\Sigma=\Sigma^-\sqcup\Sigma^+$。称 $S=\widehat S\setminus\Sigma$ 为\emph{边界带标点的 Riemann 曲面}，称 $\Sigma^-$、$\Sigma^+$ 中的点分别为它的\emph{入向无穷远点}和\emph{出向无穷远点}；这些点仍作为数据的一部分保留。注意，可以从 $S$ 典范地恢复 $\widehat S$。对于其中最简单的曲面，采用以下记号：
\begin{itemize}
\item $D$ 是 $\CC$ 中的闭单位圆盘；
\item $H$ 是闭上半平面，有一个入向无穷远点；
\item $\overline H$ 是与 $H$ 相同的曲面，但将其无穷远点视为出向的；
\item $Z=\R\times[0,1]$，坐标为 $(s,t)$，入向无穷远点为 $s=-\infty$，出向无穷远点为 $s=+\infty$。
\end{itemize}
边界带标点的 Riemann 曲面的\emph{同构}，是保持入向与出向无穷远点之区分的双全纯映射。
\end{definition*}
\begin{definition*}
$S$ 的一组 \emph{Lagrangian 标记}是以连通分支 $C\subset\partial S$ 为指标的一族正合 Lagrangian 子流形 $\{L_C\}$，其中 $L_C\subset M$。每个 $\zeta\in\Sigma\subset\widehat S$ 都属于两个边界分支 $C_{\zeta,0}$、$C_{\zeta,1}$ 的闭包；这两个分支
\SourcePage{109}
可以相同。因此，给定一组 Lagrangian 标记后，就有与之对应的一对子流形 $(L_{\zeta,0},L_{\zeta,1})$。约定：当 $\zeta\in\Sigma^-$ 时，按边界的自然定向，$C_{\zeta,1}$ 位于 $\zeta$ 之前，$C_{\zeta,0}$ 位于其后；对于 $\zeta\in\Sigma^+$，约定相反。
\end{definition*}

\subsection{Floer 上同调}\label{sec:8b}
固定一个特征为 $2$ 的域 $\K$。对于 $M$ 中的每一对正合 Lagrangian 子流形 $(L_0,L_1)$，都有一个 Floer 上同调群 $HF^{\mathrm{pr}}(L_0,L_1)$，它是有限维 $\K$-向量空间。这里的 $\mathrm{pr}$ 是“初步版本”的缩写，用来与后面将介绍的更精细的理论区分；在后者中，$\operatorname{char}(\K)$ 可以任意，而且 Floer 群带有分次。对每个带有 Lagrangian 标记的边界带标点 Riemann 曲面 $S$，都可以赋予一个同态
\begin{equation}\label{eq:8.1}
\Phi_S:\bigotimes_{\zeta^+\in\Sigma^+}HF^{\mathrm{pr}}(L_{\zeta^+,0},L_{\zeta^+,1})
\longrightarrow\bigotimes_{\zeta^-\in\Sigma^-}HF^{\mathrm{pr}}(L_{\zeta^-,0},L_{\zeta^-,1}).
\end{equation}
这个反变约定可能显得有些奇怪，但要记住，这里讨论的是 Floer\emph{上同调}。这些映射满足开弦拓扑量子场论（TQFT）公理的一种形式，参见例如\bcite{122}；其中最重要的是下面的\emph{复合律}。设 $S_1,S_2$ 是带有 Lagrangian 标记的边界带标点 Riemann 曲面。假设有无穷远点 $\zeta_1^+\in\Sigma_1^+$、$\zeta_2^-\in\Sigma_2^-$，它们配备同一\emph{有序}对子流形。像\sr{fig:8.1}{图~8.1} 那样将两者粘合，得到另一个边界带标点 Riemann 曲面 $S=S_1\#S_2$。由于 TQFT 的另一性质——拓扑不变性——具体怎样完成粘合并不重要。$S$ 带有诱导的 Lagrangian 标记，且
\begin{equation}\label{eq:8.2}
\Phi_S=(\Id\otimes\cdots\otimes\Phi_{S_1}\otimes\cdots\otimes\Id)
\circ(\Id\otimes\cdots\otimes\Phi_{S_2}\otimes\cdots\otimes\Id).
\end{equation}
其规则是：将 $\Phi_{S_2}$ 输出中的 $\zeta_2^-$ 因子，代入 $\Phi_{S_1}$ 输入中的 $\zeta_1^+$ 分量。\znote{101}{反变约定的来源是：Floer 群是路径空间上作用量泛函的 Morse 理论，而带标点的曲面 $S$ 给出的是「推前」——把入向端的闭链沿模空间推到出向端；翻译成上同调的语言，这就成了从出向因子到入向因子的映射，方向随之反转。若把 $L$ 看作对象、$HF(L_0,L_1)$ 看作态射空间，同一件事就是复合的结合律，这也是后文把 $H\F(M)$ 称作范畴、把 $\Phi_S$ 称作复合的原因。[来源：nLab, "Floer homology", https:/\allowbreak /\allowbreak ncatlab.org/\allowbreak nlab/\allowbreak show/\allowbreak Floer+\allowbreak homology（2026-11-15 访问）；并见原书 (8.2)--(8.3)。]}

\subsection{Donaldson–Fukaya 范畴}\label{sec:8c}
一般而言，任意开弦 TQFT 都对应一个 Frobenius 范畴。在 Floer 理论中，由此得到一个 $\K$-线性范畴 $H\F(M)^{\mathrm{pr}}$，称为 Donaldson–Fukaya 范畴的初步版本。对象是正合 Lagrangian 子流形 $L\subset M$，从 $L_0$ 到 $L_1$ 的态射空间是 Floer 群 $HF^{\mathrm{pr}}(L_0,L_1)$。复合
\begin{equation}\label{eq:8.3}
HF^{\mathrm{pr}}(L_1,L_2)\otimes HF^{\mathrm{pr}}(L_0,L_1)\longrightarrow HF^{\mathrm{pr}}(L_0,L_2)
\end{equation}
是\emph{三角形乘积}。用 TQFT 的语言来说，它就是不变量 $\Phi_S$：这里 $S$ 是去掉三个边界点的圆盘，其中一个点入向、两个点出向，并用 $L_0,L_1,L_2$ 标记。恒等态射是 $\Phi_H\in HF^{\mathrm{pr}}(L,L)$，其中在 $\partial H$ 上放置标记 $L$。
\SourcePage{110}
\SourceFigure{figure-8-1}{0.50469176}{8.1}
复合的结合律，以及 $\Phi_H$ 是单位元这一事实，都由 TQFT 的粘合公理推出。$H\F(M)^{\mathrm{pr}}$ 的 Frobenius 性质如下：对于每个 $L$，不变量 $\Phi_{\overline H}$ 给出一个线性映射 $HF^{\mathrm{pr}}(L,L)\to\K$。将其与复合结合，对于任意 $L_0,L_1$，得到非退化配对
\[
HF^{\mathrm{pr}}(L_1,L_0)\otimes_{\K}HF^{\mathrm{pr}}(L_0,L_1)\longrightarrow\K.
\]
更直接地，定义这一配对的曲面，是两个无穷远点都出向的无限条带。接上这样的条带，可以把入向无穷远点换成出向无穷远点。这意味着两类点的区别在一定程度上是人为的。确实，可以不作这种区别，从而使这里的公理体系与 TQFT 的原始表述\bcite{7}一致；但由于后面会说明的技术原因，我们选择保留它。\znote{102}{保留这一区分有两个实际理由：条带状端与非齐次项 $K\in\Omega^1(S,\mathcal H)$ 都要沿 $s$ 方向定义（见 (8.6)--(8.7)），而 Floer 方程 (8.8) 本身就不对称——$\partial_su+J(\partial_tu-X)=0$ 把 $s$ 与 $t$ 区别对待。若抹掉入向与出向的区别，得到的是 TQFT 的原始公理，但那时无法直接写出上链层次的映射 $C\Phi_S$；本章保留区分，正是为了能逐点定义计数：粘合两张曲面必须一进一出，若两端同类便无法直接接上。[来源：原书 §8b 本段与 §8d--8f 的构造，并见原书文献 [7]。]}

还应指出 Donaldson–Fukaya 范畴的两个更具特殊性的性质。首先，对于任意 $L$，都有一个 Frobenius 代数的典范同构，即 Piunikhin–Salamon–Schwarz（PSS）映射\bcite{106}：
\begin{equation}\label{eq:8.4}
H^*(L;\K)\cong HF^{\mathrm{pr}}(L,L).
\end{equation}
理解\eref{8.4}的最佳方式，是将它视为 TQFT 框架的一种自然推广的一部分：在这种推广中，曲面 $S$ 还带有额外的边界标点，对这些点的处理类似于 Gromov–Witten 理论。

其次，每个正合 Lagrangian 子流形的同痕 $\{L_s\}_{0\leq s\leq1}$ 都给出 $HF^{\mathrm{pr}}(L_0,L_1)$ 中的一个特殊元素，具有以下性质：它只依赖于 $\{L_s\}$ 相对于端点的同伦类；对于
\SourcePage{111}
常值同痕，它是恒等态射；同痕作为路径的复合，对应于取相关元素的三角形乘积。直接推论是：这些元素在 $H\F(M)^{\mathrm{pr}}$ 中全都是同构。因此，同痕的 Lagrangian 子流形给出 Donaldson–Fukaya 范畴中同构的对象，进而 Floer 上同调群具有同痕不变性。从形式上说，这一构造以基本 TQFT 结构的另一种扩张为基础；其中允许 Lagrangian 边界条件沿 $\partial S$ 上的点光滑变化。

\subsection{条带状端}\label{sec:8d}
现在开始回顾 Floer 理论的基本基础。正如本章开头列举的参考文献所表明的，这些都是熟悉的材料；但对于后面讨论 Fukaya 范畴，它们是很好的预备练习。

首先，虽然边界带标点的 Riemann 曲面已经足以表述 TQFT，但在实际构造时，采用稍为刚性的结构更好。
\begin{definition*}
将半无限条带记为 $Z^\pm=\R^\pm\times[0,1]$。设 $S$ 是边界带标点的 Riemann 曲面，无穷远点集合为 $\Sigma^\pm$。$S$ 的一组\emph{条带状端}，是对每个 $\zeta\in\Sigma^\pm$ 给定的固有全纯嵌入 $\epsilon_\zeta:Z^\pm\to S$，满足
\[
\epsilon_\zeta^{-1}(\partial S)=\R^\pm\times\{0,1\},\qquad
\lim_{s\to\pm\infty}\epsilon_\zeta(s,\cdot)=\zeta,
\]
并且各 $\epsilon_\zeta$ 的像两两不交。$Z^\pm$ 上的 Riemann 曲面结构可延伸到单点紧化 $\widehat Z^\pm=Z^\pm\cup\{\pm\infty\}$，每个条带状端也都延伸为全纯嵌入 $\widehat\epsilon_\zeta:\widehat Z^\pm\to\widehat S$，将 $\pm\infty$ 映到 $\zeta$。
\end{definition*}
需要这些端有几个原因。其一是，作为理论的分析框架所用的 $W^{1,p}$（$p>2$）范数不具有共形不变性。\znote{103}{选 $p>2$ 的技术原因：二维 Sobolev 嵌入 $W^{1,p}\subset C^0$ 在 $p>2$ 时成立，故 $u$ 连续，边界条件 $u(\partial S)\subset L_C$ 才有逐点意义；$p=2$ 时须改用迹映射陈述（见 §8h）。另一方面 $p=2$ 恰是方程共形不变的情形，此时 $W^{1,2}$ 范数不具共形不变性，模空间在共形变动下不易控制，故实际取 $p>2$，并辅以条带状端上的标准坐标把非齐次项搬过去。[来源：原书 §8d 本段与 §8h 关于 $W^{1,p}$、迹映射的讨论，并见原书文献 [29]。]}同样重要的是，定义 Floer 上同调时出现的非齐次项涉及 $Z=\R\times[0,1]$ 上的标准坐标，因此为了定义运算\eref{8.1}，需要在更一般的 Riemann 曲面的端上赋予类似结构。还有一点是，可以使\sr{fig:8.1}{图~8.1} 的粘合过程更加精确：给定 $\zeta_1^+$、$\zeta_2^-$ 和粘合长度 $l>0$，令
\begin{equation}\label{eq:8.5}
\begin{aligned}
S'_1&=S_1\setminus\epsilon_{\zeta_1^+}((l,\infty)\times[0,1]),\\
S'_2&=S_2\setminus\epsilon_{\zeta_2^-}((-\infty,-l)\times[0,1]),\\
S=S_1\#_l S_2&=(S'_1\sqcup S'_2)/\sim,
\end{aligned}
\end{equation}
其中对 $(s,t)\in[0,l]\times[0,1]$，作识别 $\epsilon_{\zeta_2^-}(s-l,t)\sim\epsilon_{\zeta_1^+}(s,t)$。

\subsection{Floer 数据与扰动数据}\label{sec:8e}
将 $M$ 上所有与 $\varpi$ 相容、且在边界附近与给定的 $I$ 相同的殆复结构 $J$ 所成的空间记为 $\mathcal J$。若 $T$ 是任意流形，则用 $C^\infty(T,\mathcal J)$ 表示由 $T$ 光滑参数化的殆复结构族
\SourcePage{112}
$J=\{J(t)\}_{t\in T}$ 所成的空间，要求它们在 $T\times U$ 上与 $I$ 相同，其中 $U$ 是 $\partial M$ 的某个邻域。类似地，用 $\mathcal H=C_c^\infty(\operatorname{int}(M),\R)$ 表示 $M$ 上在边界附近为零的光滑函数所成的空间；$C^\infty(T,\mathcal H)$ 表示 $T\times M$ 上在某个如上的 $T\times U$ 上为零的函数。
\begin{definition*}
设 $L_0,L_1\subset M$ 是一对正合 Lagrangian 子流形。\emph{Floer 数据}由 $H\in C^\infty([0,1],\mathcal H)$、$J\in C^\infty([0,1],\mathcal J)$ 组成，满足以下性质：若 $X$ 是 $H$ 的时变 Hamilton 向量场，$\phi$ 是它的“流”，则 $\phi^1(L_0)$ 与 $L_1$ 横截相交。
\end{definition*}
设 $S$ 是带有 Lagrangian 标记的边界带标点 Riemann 曲面。假设已经为它选定条带状端，并且对每个无穷远点 $\zeta\in\Sigma$ 所对应的子流形对 $(L_{\zeta,0},L_{\zeta,1})$ 选定 Floer 数据 $(H_\zeta,J_\zeta)$。
\begin{definition*}
$S$ 的\emph{扰动数据}是一对 $(K,J)$，其中 $K\in\Omega^1(S,\mathcal H)$ 满足
\begin{equation}\label{eq:8.6}
K(\xi)|_{L_C}=0\qquad\text{对所有 }\xi\in TC\subset T(\partial S).
\end{equation}
换言之，$K$ 是取值于 $M$ 上函数空间的一形式；如果取一个属于连通分支 $C\subset\partial S$ 的边界点，以及该点处与边界相切的向量 $\xi$，那么 $K(\xi)\in\mathcal H$ 是沿 $L_C\subset M$ 为零的函数。扰动数据的第二部分，就是殆复结构族 $J\in C^\infty(S,\mathcal J)$。此外，$K,J$ 必须与选定的条带状端及 Floer 数据相容，即
\begin{equation}\label{eq:8.7}
\epsilon_\zeta^*K=H_\zeta(t)\,dt,\qquad J(\epsilon_\zeta(s,t))=J_\zeta(t)
\end{equation}
对每个 $\zeta\in\Sigma^\pm$ 和 $(s,t)\in Z^\pm$ 成立。
\end{definition*}
注意，反过来，在 $S=Z$ 上令 $K(s,t)=H(t)\,dt$，便可将 Floer 数据视为平移不变的扰动数据。Floer 数据与扰动数据这两个概念，规定了要加入伪全纯映射方程中的非齐次项的类型。

\subsection{非齐次伪全纯映射}\label{sec:8f}
设有一对正合 Lagrangian 子流形 $L_0,L_1$，以及 Floer 数据 $(H,J)$。像前面一样，以 $X$ 表示 $H$ 的时变向量场。定义 $\mathcal C(L_0,L_1)$ 为满足 $y(0)\in L_0$、$y(1)\in L_1$ 及 $dy/dt=X(t,y(t))$ 的映射 $y:[0,1]\to M$ 所成的集合。它们与 $\phi^1(L_0)\cap L_1$ 中的点一一对应，因此只有有限个。对 $u\in C^\infty(Z,M)$，Floer 方程为
\begin{equation}\label{eq:8.8}
\begin{cases}
\partial_su+J(t,u)(\partial_tu-X(t,u))=0,\\
u(s,0)\in L_0,\quad u(s,1)\in L_1.
\end{cases}
\end{equation}
以 $\M_Z(y_0,y_1)$ 表示\eref{8.8}满足渐近条件 $\lim_{s\to+\infty}u(s,\cdot)=y_1$、$\lim_{s\to-\infty}u(s,\cdot)=y_0$ 的解的集合，其中 $y_0,y_1\in\mathcal C(L_0,L_1)$。规定这种收敛是 $C^\infty$ 意义下的指数速率收敛；由于横截
\SourcePage{113}
相交假设，所有合理的收敛条件都是等价的。通过 $s$ 变量的平移，$\R$ 作用在 $\M_Z(y_0,y_1)$ 上。当 $y_0\ne y_1$ 时，作用是自由的，此时将商空间记为 $\M_Z^*(y_0,y_1)$。在 $y_0=y_1$ 的例外情形，$\M_Z(y_0,y_1)$ 是一个点，唯一的元素是驻定解 $u(s,t)=y_0(t)$；此时约定 $\M_Z^*(y_0,y_1)=\varnothing$。

更一般地，取一个带有 Lagrangian 标记的边界带标点 Riemann 曲面 $S$，为它配备条带状端、每个无穷远点处的 Floer 数据 $(H_\zeta,J_\zeta)$，以及相容的扰动数据 $(K,J)$。$K$ 确定一个向量场值一形式 $Y\in\Omega^1(S,C^\infty(TM))$：对于每个 $\xi\in TS$，$Y(\xi)$ 是 $K(\xi)$ 的 Hamilton 向量场。对 $u\in C^\infty(S,M)$，\emph{非齐次伪全纯映射方程}为
\begin{equation}\label{eq:8.9}
\begin{cases}
Du(z)+J(z,u)\circ Du(z)\circ I_S
=Y(z,u)+J(z,u)\circ Y(z,u)\circ I_S,\\
u(C)\subset L_C\quad\text{对所有 }C\subset\partial S,
\end{cases}
\end{equation}
其中 $I_S$ 是 $S$ 的复结构。可以将它更简洁地写成 $(Du-Y)^{0,1}=0$，或写成 $\overline\partial_J(u)=\nu(z,u(z))$，其中 $\nu=Y^{0,1}$。与\bcite{115}比较可知，我们没有使用最一般的非齐次项 $\nu$。显然，Floer 方程是\eref{8.9}的特例，其中 $S=Z$，数据 $(K,J)$ 平移不变。反过来，\eref{8.7}意味着\eref{8.9}在 $S$ 的条带状端上化为\eref{8.8}。因此，可以考虑满足以下收敛条件的\eref{8.9}的解集 $\M_S(\{y_\zeta\})$：
\begin{equation}\label{eq:8.10}
\lim_{s\to\pm\infty}u(\epsilon_\zeta(s,\cdot))=y_\zeta,
\qquad \zeta\in\Sigma^\pm,\quad y_\zeta\in\mathcal C(L_{\zeta,0},L_{\zeta,1}).
\end{equation}

\subsection{能量}\label{sec:8g}
先考虑 $L_0,L_1$ 及 Floer 数据 $(H,J)$。由于正合性，存在函数 $h_{L_k}\in C^\infty(L_k,\R)$，满足 $dh_{L_k}=\theta|_{L_k}$。路径空间 $\{y\in C^\infty([0,1],M):y(0)\in L_0,\ y(1)\in L_1\}$ 上的 $H$-扰动作用量泛函为
\begin{equation}\label{eq:8.11}
A_H(y)=\left(\int -y^*\theta+H(t,y(t))\,dt\right)+h_{L_1}(y(1))-h_{L_0}(y(0)).
\end{equation}
$A_H$ 的临界点集恰好是 $\mathcal C(L_0,L_1)$，\eref{8.8}则是某个 $L^2$ 度量下的负梯度流方程。\znote{105}{(8.8) 是 $A_H$ 在某个 $L^2$ 度量下的负梯度流方程，这一点值得记住：$A_H$ 的临界点集正是 $\mathcal C(L_0,L_1)$，而 (8.8) 正是从 $y_0$ 流向 $y_1$ 的下降轨道。于是 $E(u)=\int_Z|\partial_su|^2=A_H(y_0)-A_H(y_1)$ 就是 Morse 理论中「能量等于势差」的无穷维版本；它把解的能量换成端点数据，紧性论证才有统一的能量界可用。这也正是 (8.12) 的实质：先验界不是估计出来的，而是由端点决定，这也正是它被称为先验界的理由。[来源：nLab, "Floer homology", https:/\allowbreak /\allowbreak ncatlab.org/\allowbreak nlab/\allowbreak show/\allowbreak Floer+\allowbreak homology（2026-11-15 访问）；并见原书 (8.11)--(8.12)。]}因此，对任意 $u\in\M_Z(y_0,y_1)$，都有先验界
\[
E(u)=\int_Z|\partial u/\partial s|^2=A_H(y_0)-A_H(y_1).
\]
为了将这一点推广到\eref{8.9}，对每个连通分支 $C\subset\partial S$ 选取具有同样性质的函数 $h_{L_C}\in C^\infty(L_C,\R)$。于是，对 $u\in\M_S(\{y_\zeta\})$，有
\begin{equation}\label{eq:8.12}
E(u)\overset{\mathrm{def}}=\int_S\tfrac12|Du-Y|^2
=\sum_{\zeta\in\Sigma^\pm}\mp A_{H_\zeta}(y_\zeta)-\int_S R(z,u(z)).
\end{equation}
在对 $\zeta$ 求和时，每个 $h_{L_C}$ 的选择所含的常数歧义——不连通时为局部常数歧义——都会抵消，因此第一项是有意义的。
\SourcePage{114}
至于第二项，曲率被积形式 $R\in\Omega^2(S,\mathcal H)$ 在 $S$ 的局部坐标 $(s,t)$ 中为
\[
R=\bigl(\partial_sK(\partial_t)-\partial_tK(\partial_s)
+\{K(\partial_s),K(\partial_t)\}\bigr)\,ds\wedge dt.
\]
这里 $\{\cdot,\cdot\}$ 是 Poisson 括号。由于 $R$ 在 $S$ 的端上为零，积分有一个与映射 $u$ 无关的界，因此 $u\in\M_S(\{y_\zeta\})$ 的能量 $E(u)$ 也有这样的界。

\subsection{Cauchy–Riemann 算子}\label{sec:8h}
设 $S$ 是带边 Riemann 曲面，$E\to S$ 是带有辛结构 $\varpi_E$ 和相容复结构 $J_E$ 的向量丛，$F\subset E|_{\partial S}$ 是 Lagrangian 子丛。用 $C^\infty(S,E,F)$ 表示在边界上的限制取值于 $F$ 的截面空间，用 $C^\infty(S,\Omega_S^{0,1}\otimes E)$ 表示 $E$-值 $(0,1)$-形式的空间，不附加任何边界条件。对于较低的光滑性，同样可以定义局部 $W^{1,p}$ 截面的空间 $W_{\mathrm{loc}}^{1,p}(S,E,F)$，其中 $2\leq p<\infty$。当 $p>2$ 时，截面连续，边界条件显然有意义；当 $p=2$ 时，通过限制映射 $W_{\mathrm{loc}}^{1,2}(S,E)\to L_{\mathrm{loc}}^2(\partial S,E)$ 定义边界条件，参见例如\bcite{29}。还可定义局部 $L^p$ 的 $(0,1)$-形式空间 $L_{\mathrm{loc}}^p(S,\Omega_S^{0,1}\otimes E)$。给定 $E$ 上任意联络 $\nabla$，不必是复联络或辛联络，将对应的 Cauchy–Riemann 算子定义为协变微分的 $(0,1)$ 部分：
\begin{equation}\label{eq:8.13}
\overline\partial_\nabla=\nabla^{0,1}:C^\infty(S,E,F)\longrightarrow C^\infty(S,\Omega_S^{0,1}\otimes E).
\end{equation}
它可延伸为 $W_{\mathrm{loc}}^{1,p}(S,E,F)\to L_{\mathrm{loc}}^p(S,\Omega_S^{0,1}\otimes E)$。若 $S$ 紧致，则这一延伸是 Fredholm 算子。

为了在非紧致曲面上建立\eref{8.13}的合理分析理论，需要假设 $S$ 带有条带状端，并且 $(E,F,\nabla)$ 在这些端上的渐近行为受到充分控制。精确表述需要一些预备定义。
\begin{definition*}
对曲面 $S$ 的每个端 $\zeta$，给定向量丛 $E_\zeta\to[0,1]$，配备辛结构 $\varpi_{E_\zeta}$、相容复结构 $J_{E_\zeta}$、\emph{辛联络} $\nabla_\zeta$，以及 $k=0,1$ 时的 Lagrangian 子空间 $\Lambda_{\zeta,k}\subset E_{\zeta,k}$。称 $(E_\zeta,\varpi_{E_\zeta},J_{E_\zeta},\nabla_\zeta,\Lambda_{\zeta,0},\Lambda_{\zeta,1})$ 为 $\zeta$ 处的\emph{极限数据}。通常方便假设这一数据\emph{非退化}，意即沿 $[0,1]$ 的 $\nabla_\zeta$-平行移动把 $\Lambda_{\zeta,0}$ 映到与 $\Lambda_{\zeta,1}$ 横截的子空间。
\end{definition*}
设 $\pi^\pm:Z^\pm\to[0,1]$ 是沿 $t$ 方向的投影。要求向量丛 $E$ 配备指定的同构
\begin{equation}\label{eq:8.14}
\epsilon_\zeta^*E\cong(\pi^\pm)^*E_\zeta,
\end{equation}
并且这些同构在渐近意义下与极限数据相容。也就是说，用\eref{8.14}识别两端后，差 $\varpi_E-\varpi_{E_\zeta}$、$J_E-J_{E_\zeta}$、$\nabla-\nabla_\zeta$ 在 $s\to\pm\infty$ 时以指数速率衰减，例如在 $C^1$ 拓扑中；同时，$F_{\epsilon_\zeta(s,k)}\to\Lambda_{\zeta,k}$ 也在同样意义下成立，这里的收敛在 $(E_\zeta)_k$ 的 Lagrangian 子空间 Grassmann 流形中理解。
\SourcePage{115}
\begin{definition*}
若 $\overline\partial_\nabla$ 在各端满足\eref{8.14}及上述渐近条件，就称其为\emph{容许的}。若所有极限数据还都是非退化的，就称 $\overline\partial_\nabla$ 本身\emph{非退化}。
\end{definition*}
端上的平凡化确定一个沿 $s$ 方向不变的度量，用它来定义\emph{整体} $W^{1,p}$ 截面及 $L^p$ 截面的空间 $W^{1,p}(S,E,F)$、$L^p(S,\Omega_S^{0,1}\otimes E)$。于是\eref{8.13}延伸为连续算子
\begin{equation}\label{eq:8.15}
\overline\partial_\nabla:W^{1,p}(S,E,F)\longrightarrow L^p(S,\Omega_S^{0,1}\otimes E).
\end{equation}
若满足前述非退化条件，这些延伸都是 Fredholm 算子，其核与余核不依赖于 $p$。更一般的做法，是引入在各端带有某个权重 $\epsilon\geq0$ 的空间 $W^{1,p;\epsilon}$、$L^{p;\epsilon}$；这些空间的元素是指数衰减的截面，当 $\epsilon=0$ 时则退化为先前考虑的通常 Sobolev 空间。对于任意容许的 $\overline\partial$ 算子，都存在 $\epsilon_0>0$，使得当 $\epsilon\in(0,\epsilon_0)$ 时，延伸
\begin{equation}\label{eq:8.16}
\overline\partial_\nabla:W^{1,p;\epsilon}(S,E,F)\longrightarrow L^{p;\epsilon}(S,\Omega_S^{0,1}\otimes E)
\end{equation}
是 Fredholm 算子，其核与余核不依赖于权重，比较\bcite{45}，命题 4.1。若 Cauchy–Riemann 算子非退化，同一结论对 $\epsilon\in[0,\epsilon_0)$ 也成立；这意味着在此情形，\eref{8.16}的核与余核分别与\eref{8.15}相同。

有时还将使用带有更高阶导数 $k>1$ 的 Sobolev 空间 $W^{k,p}$，以及加权版本 $W^{k,p;\epsilon}$；上述讨论同样适用于它们。
